\documentclass[12pt,a4paper,reqno]{amsart}

% language
\usepackage[greek,english]{babel}
\usepackage[utf8]{inputenc}
\usepackage{alphabeta}

% change default names to greek
\addto\captionsenglish{
    \renewcommand{\contentsname}{Περιεχόμενα}
    \renewcommand{\refname}{Βιβλιογραφία}
    \renewcommand{\datename}{Ημερομηνία:}
    \renewcommand{\urladdrname}{Ιστοσελίδα}
}

% math 
\usepackage{amsmath,amsthm,amssymb,amscd}

% font
\usepackage[cal=euler]{mathalfa}
\usepackage{libertinus-type1}
% \usepackage{txfonts} % for upright greek letters
\usepackage{bm} % for bold symbols
\usepackage{bbm} % for the simply-looking bb symbols

% miscellaneous 
\usepackage{changepage} %for indenting environments
\usepackage{csquotes} % example: \textcquote{}
\usepackage{blkarray}
\setcounter{MaxMatrixCols}{20} % default for pmatrix is 10!!
\usepackage{ytableau}
\usepackage{array} %needed to increase the vertical length in a tabular

% drawing
\usepackage{tikz,tikz-cd}
\usetikzlibrary{shapes.misc, patterns, matrix, calc, intersections,positioning,arrows.meta}
\usepackage{graphics,graphicx}
\usepackage{float} % provides enhanced control and customization options for floating objects such as figures and tables

% colors
\usepackage{xcolor}
\definecolor{darkcandyapplered}{rgb}{0.64, 0.0, 0.0}
\definecolor{midnightblue}{rgb}{0.1, 0.1, 0.44}
\definecolor{mylightblue}{HTML}{336699}
\definecolor{burntorange}{rgb}{0.8, 0.33, 0.0}
\definecolor{iceberg}{rgb}{0.44, 0.65, 0.82}
\definecolor{applegreen}{rgb}{0.55, 0.71, 0.0}
\definecolor{canaryyellow}{rgb}{1.0, 0.94, 0.0}
\definecolor{lightgray}{rgb}{0.83, 0.83, 0.83}

% hrefs
\usepackage{hyperref}
\usepackage[noabbrev,capitalize]{cleveref}
\hypersetup{
    pdftoolbar=true,        
    pdfmenubar=true,        
    pdffitwindow=false,     
    pdfstartview={FitH},  % fits the width of the page to the window
    pdftitle={},
    pdfauthor={},
    pdfsubject={},
    pdfkeywords={},
    pdfnewwindow=true,  % links in new window
    colorlinks=true,  % false: boxed links; true: colored links
    linkcolor=darkcandyapplered,   % color of internal links
    citecolor=midnightblue,  % color of links to bibliography
    urlcolor=cyan,  % color of external links
    linktocpage=true  % changes the links from the section body to the page number
    }

% geometry
\textwidth=16cm 
\textheight=21cm 
\hoffset=-55pt 
\footskip=25pt

% thm envs (you might need to change the path)
\theoremstyle{definition}
\newtheorem{exercise}{Άσκηση}
\newtheorem{solution}{Λύση}
\newtheorem*{remark}{Παρατήρηση}
\newtheorem*{example}{Παράδειγμα}

% math ops 
% fields
\newcommand{\NN}{\mathbbmss{N}} 
\newcommand{\ZZ}{\mathbbmss{Z}} 
\newcommand{\QQ}{\mathbbmss{Q}} 
\newcommand{\RR}{\mathbbmss{R}} 
\newcommand{\CC}{\mathbbmss{C}} 
\newcommand{\KK}{\mathbbmss{K}} 
\newcommand{\FF}{\mathbbmss{F}} 
\newcommand{\HH}{\mathbbmss{H}} 

% symmetric group
\newcommand{\fS}{\mathfrak{S}}  
\newcommand{\fp}{\mathfrak{p}}  
\newcommand{\fm}{\mathfrak{m}}  

% calligraphic 
\newcommand{\aA}{\mathcal{A}} 
\newcommand{\bB}{\mathcal{B}}
\newcommand{\cC}{\mathcal{C}}
\newcommand{\dD}{\mathcal{D}}
\newcommand{\eE}{\mathcal{E}}
\newcommand{\fF}{\mathcal{F}}
\newcommand{\hH}{\mathcal{H}}
\newcommand{\iI}{\mathcal{I}}
\newcommand{\lL}{\mathcal{L}}
\newcommand{\nN}{\mathcal{N}}
\newcommand{\oO}{\mathcal{O}}
\newcommand{\pP}{\mathcal{P}}
\newcommand{\sS}{\mathcal{S}}
\newcommand{\mM}{\mathcal{M}}
\newcommand{\uU}{\mathcal{U}}

% bold
\newcommand{\bfa}{\mathbf{a}}
\newcommand{\bfe}{\mathbf{e}}
\newcommand{\bfF}{\pmb{F}}
\newcommand{\bfR}{\pmb{R}}
\newcommand{\bfv}{\mathbf{v}}
%\newcommand{\bfx}{\bm{x}}
%\newcommand{\bfx}{\mathbf{x}} 
\newcommand{\bfx}{\pmb{x}}
\newcommand{\bfX}{\pmb{X}}
\newcommand{\bfy}{\pmb{y}}
\newcommand{\bfz}{\pmb{z}}

% roman
\newcommand{\rmA}{\mathrm{A}}
\newcommand{\rmB}{\mathrm{B}}
\newcommand{\rmC}{\mathrm{C}}
\newcommand{\rmD}{\mathrm{D}} 
\newcommand{\rmH}{\mathrm{H}} 
\newcommand{\rmh}{\mathrm{h}} 
\newcommand{\rmI}{\mathrm{I}} 
\newcommand{\rmJ}{\mathrm{J}} 
\newcommand{\rmK}{\mathrm{K}}
\newcommand{\rmM}{\mathrm{M}}
\newcommand{\rmP}{\mathrm{P}}  
\newcommand{\rmp}{\mathrm{p}}  
\newcommand{\rmQ}{\mathrm{Q}}  
\newcommand{\rmR}{\mathrm{R}}
\newcommand{\rmS}{\mathrm{S}}
\newcommand{\rmT}{\mathrm{T}}
\newcommand{\rmU}{\mathrm{U}}
\newcommand{\rmV}{\mathrm{V}}
\newcommand{\rmY}{\mathrm{Y}}
\newcommand{\rmZ}{\mathrm{Z}}
\newcommand{\rmz}{\mathrm{z}}

% arrows and symbols 
\renewcommand{\to}{\rightarrow}
\newcommand{\toto}{\longrightarrow}
\newcommand{\mapstoto}{\longmapsto}
\newcommand{\then}{\Rightarrow}
\newcommand{\IFF}{\Leftrightarrow}
\newcommand{\tl}{\tilde}
\newcommand{\wtl}{\widetilde}
\newcommand{\ol}{\overline}
\newcommand{\ul}{\underline}
\newcommand{\oldemptyset}{\emptyset}
\renewcommand{\emptyset}{\varnothing}
\DeclareMathSymbol{\Arg}{\mathbin}{AMSa}{"39} % for arguments 
\newcommand{\onto}{\ensuremath{\twoheadrightarrow}}
\newcommand{\tle}{\trianglelefteq}
\newcommand{\tge}{\trianglerighteq}

% custom ops
\newcommand{\rmKer}{\mathrm{Ker}}
\newcommand{\rmIm}{\mathrm{Im}}
\newcommand{\lcm}{\mathrm{lcm}}
\newcommand{\GL}{\mathrm{GL}}
\newcommand{\ev}{\mathrm{ev}}
\newcommand{\End}{\mathrm{End}}
\newcommand{\rmid}{\mathrm{id}}
\newcommand{\rmspan}{\mathrm{span}}
\newcommand{\Spec}{\mathrm{Spec}}
\newcommand{\mSpec}{\mathrm{mSpec}}
\newcommand{\rad}{\mathrm{rad}}
\newcommand{\Rad}{\mathrm{Rad}}
\newcommand{\Jac}{\mathrm{Jac}}
\newcommand{\tor}{\mathrm{tor}}
\newcommand{\Ann}{\mathrm{Ann}}
\newcommand{\Hom}{\mathrm{Hom}}

% absolute value symbol
\usepackage{mathtools} 
\DeclarePairedDelimiter\abs{\lvert}{\rvert}%
\DeclarePairedDelimiter\norm{\lVert}{\rVert}%
\makeatletter
\let\oldabs\abs
\def\abs{\@ifstar{\oldabs}{\oldabs*}}

% tensor symbol
\newcommand{\tensor}[1]{%
  \mathbin{\mathop{\otimes}\limits_{#1}}%
}

% permutation cycle notation
\ExplSyntaxOn
\NewDocumentCommand{\cycle}{ O{\;} m }
 {
  (
  \alec_cycle:nn { #1 } { #2 }
  )
 }

\seq_new:N \l_alec_cycle_seq
\cs_new_protected:Npn \alec_cycle:nn #1 #2
 {
  \seq_set_split:Nnn \l_alec_cycle_seq { , } { #2 }
  \seq_use:Nn \l_alec_cycle_seq { #1 }
 }
\ExplSyntaxOff

% setminus symbol
\newcommand{\mysetminusD}{\hbox{\tikz{\draw[line width=0.6pt,line cap=round] (3pt,0) -- (0,6pt);}}}
\newcommand{\mysetminusT}{\mysetminusD}
\newcommand{\mysetminusS}{\hbox{\tikz{\draw[line width=0.45pt,line cap=round] (2pt,0) -- (0,4pt);}}}
\newcommand{\mysetminusSS}{\hbox{\tikz{\draw[line width=0.4pt,line cap=round] (1.5pt,0) -- (0,3pt);}}}
\newcommand{\sm}{\mathbin{\mathchoice{\mysetminusD}{\mysetminusT}{\mysetminusS}{\mysetminusSS}}}

% miscellaneous commands
\newcommand{\defn}[1]{{\color{mylightblue}{#1}}}
\newcommand{\toDo}{{\bf\color{red} TODO}}
\newcommand{\toCite}{{\bf\color{green} CITE}}
\newcommand*{\vertbar}{\rule[-1ex]{0.5pt}{2.5ex}} % for matrices with column vectors
\newcommand*{\horzbar}{\rule[.5ex]{2.5ex}{0.5pt}} % for matrices with row vectors
\newcommand{\myblue}[1]{{\color{iceberg}{#1}}}
\newcommand{\myorange}[1]{{\color{burntorange}{#1}}}
\newcommand{\mygreen}[1]{{\color{applegreen}{#1}}}
\newcommand{\myred}[1]{{\color{darkcandyapplered}{#1}}}

% ferrer's diagram
\newcommand{\fdiagram}[1]{
    \begin{tikzpicture}[scale=.7]
        \fill foreach \Z [count=\Y] in {#1}
        {foreach \X in {1,...,\Z} 
        {(\X,-\Y) circle[radius=3pt]}};
    \end{tikzpicture}
}

%
\usepackage{pgfplots}
\pgfplotsset{compat=1.18}

%
\makeatletter
\def\mathcolor#1#{\@mathcolor{#1}}
\def\@mathcolor#1#2#3{%
  \protect\leavevmode
  \begingroup
    \color#1{#2}#3%
  \endgroup
}
\makeatother

% 
\newenvironment{nouppercase}{%
  \let\uppercase\relax%
  \renewcommand{\uppercasenonmath}[1]{}}{}

% titlepage
\title{226: Θεωρία Δακτυλίων και Modules \\ Ασκήσεις με Ενδεικτικές Λύσεις}
\author[Β.~Δ. Μουστακας]{Βασίλης Διονύσης Μουστάκας \\ Πανεπιστήμιο Κρήτης}
\date{1 Απριλίου 2026}
% \urladdr{\href{https://sites.google.com/view/vasmous}{https://sites.google.com/view/vasmous}}

\begin{document}

\begingroup
\def\uppercasenonmath#1{} % this disables uppercase title
\let\MakeUppercase\relax % this disables uppercase authors
\maketitle
\endgroup

\setcounter{exercise}{25}
% \setcounter{theorem}{5}
\setcounter{equation}{4}
\renewcommand{\thesubsection}{\thesection.\arabic{subsection}}

Σε ότι ακολουθεί υποθέτουμε ότι $R$ είναι ένας μεταθετικός δακτύλιος και $n$ είναι ένας θετικός ακέραιος, εκτός αν αναφέρεται διαφορετικά.

\begin{exercise}{\rm(\cite[Άσκηση~10.3.12]{DF04})}
  \label{ex:DF_10.3.12}
  Αν $A, B$ και $M$ είναι $R$-πρότυπα, τότε να δείξετε ότι 
  \begin{align}
    \label{eq:10.3.12_a}
    \Hom_R(A\times{B},M) &\cong \Hom_R(A,M) \times \Hom_R(B,M)  \\
    \label{eq:10.3.12_b}
    \Hom_R(M,A\times{B}) &\cong \Hom_R(M,A) \times \Hom_R(M,B).
  \end{align}
\end{exercise}

\begin{proof}[Λύση]
  Για την Ταυτότητα \eqref{eq:10.3.12_a}, έστω $\iota_A : A \to A\times{B}$ και $\iota_B : B \to A\times{B}$, οι \textquote{φυσικοί} $R$-μονομορφισμοί, που ορίζονται θέτοντας 
  \[
  \iota_A(a) = (a,0_B) \quad \text{και} \quad \iota_B(b) = (0_A,b),
  \]
  για κάθε $a \in A$ και $b \in B$. Θεωρούμε την απεικόνιση 
  \begin{align*}
    \Phi : \Hom_R(A\times{B},M) &\to \Hom_R(A,M) \times \Hom_R(B,M) \\
    f &\mapsto (f \circ \iota_A, f \circ \iota_B). 
  \end{align*}
  Αυτή είναι $R$-γραμμική (γιατί;).
  
  Από την άλλη, θεωρούμε την απεικόνιση 
  \begin{align*}
    \Psi : \Hom_R(A,M) \times \Hom_R(B,M) &\to \Hom_R(A\times{B},M)  \\
    (f,g) &\mapsto \lambda_{f,g}, 
  \end{align*}
  όπου $\lambda_{f,g}(a,b) = f(a) + g(b)$, για κάθε $a \in A$ και $b \in B$. Η $\Psi$ είναι και αυτή $R$-γραμμική (γιατί;) και 
  \begin{equation}
    \label{eq:eq:10.3.12_a_help}
    \Phi\circ\Psi = \rmid_{\Hom_R(A,M) \times \Hom_R(B,M)}, \quad \text{και} \quad \Psi\circ\Phi = \rmid_{\Hom_R(A\times{B},M)}.
  \end{equation}
  Πράγματι, για κάθε $f \in \Hom_R(A,M)$ και $g \in \Hom_R(B,M)$, έχουμε 
  \[
    \Phi(\Psi(f,g)) = \Phi(\lambda_{f,g}) = (\lambda_{f,g}\circ\iota_A, \lambda_{f,g}\circ\iota_B).
  \]
  Όμως, 
  \[
  \lambda_{f,g}\left(
    \iota_A(a)
  \right) = \lambda_{f,g}(a,0_B) = f(a) + g(0_B) = f(a)
  \]
  για κάθε $a \in A$. Ομοίως, $\lambda_{f,g}\left(
    \iota_B(b)
  \right) = g(b)$, για κάθε $b \in B$. Επομένως, 
  \[
  \Phi(\Psi(f,g)) = (\lambda_{f,g}\circ\iota_A, \lambda_{f,g}\circ\iota_B) = (f,g)
  \]
  και έτσι έπεται η πρώτη από τις Ταυτότητες \eqref{eq:eq:10.3.12_a_help}. Όμοια έπεται και η δεύτερη και κατά συνέπεια, το ζητούμενο (γιατί;).

  Η απόδειξη της Ταυτότητας \eqref{eq:10.3.12_b} είναι ανάλογη και γι αυτό τον λόγο αφήνουμε τις λεπτομέρειες ως άσκηση. Έστω $\pi_A : A\times{B} \to A $ και $\pi_B : A\times{B} \to B$, οι \textquote{φυσικοί} $R$-επιμορφισμοί, που ορίζονται θέτοντας 
  \[
  \pi_A(a,b) = a \quad \text{και} \quad \pi_B(a,b) = b,
  \]
  για κάθε $(a,b) \in A\times B$. Θεωρούμε τις απεικονίσεις
  \begin{align*}
    \Phi : \Hom_R(M,A\times{B}) &\cong \Hom_R(M,A) \times \Hom_R(M,B) \\
    f &\mapsto (\pi_A\circ f, \pi_B\circ f) 
  \end{align*}
  και
  \begin{align*}
    \Psi : \Hom_R(M,A) \times \Hom_R(M,B) &\cong \Hom_R(M,A\times{B})  \\
    (f,g) &\mapsto \mu_{f,g}, 
  \end{align*}
  όπου $\mu_{f,g}(m) = \left(f(m), g(m)\right)$, για κάθε $m \in M$. Αυτές είναι $R$-ομομορφισμοί και αντίστροφοι η μία της άλλης.
\end{proof}

\begin{remark}
  Η Άσκηση \ref{ex:DF_10.3.12} μας πληροφορεί ότι αν ξεκινήσουμε με την βραχεία ακριβή ακολουθία 
  \[
    \begin{tikzcd}
        0 \arrow[r] & A \arrow[r, hook, "\iota_A"] & A \times B \arrow[r, two heads, "\pi_B"] & B \arrow[r] & 0
    \end{tikzcd}
  \]
  και \textquote{εφαρμόσουμε} τα $\Hom_R(-,M)$ και $\Hom(M,-)$, αντίστοιχα, προκύπτουν οι βραχείες ακριβείς ακολουθίες
  \[
    \begin{tikzcd}
        0 \arrow[r] & \Hom_R(B,M) \arrow[r, "\pi_B^\ast"] & \Hom_R(A \times B,M) \arrow[r, "\iota_A^\ast"] & \Hom_R(A,M) \arrow[r] & 0,
    \end{tikzcd}
  \]
  και
  \[
    \begin{tikzcd}
        0 \arrow[r] & \Hom_R(M,B) \arrow[r, "\prescript{\ast}{}{\pi_B}"] & \Hom_R(M, A \times B) \arrow[r, "\prescript{\ast}{}{\iota_A}"] & \Hom_R(M,A) \arrow[r] & 0,
    \end{tikzcd}
  \]
  όπου 
  \[
  \iota_A^\ast(f) = f \circ \iota_A \quad \text{και} \quad \pi_B^\ast(f) = f \circ \pi_B 
  \]
  για κάθε $f \in \Hom_R(A \times B,M)$ και 
  και 
  \[
  \prescript{\ast}{}{\iota_A}(f) = \iota_A \circ f \quad \text{και} \quad \prescript{\ast}{}{\pi_B}(f) = \pi_B \circ f,
  \]
  για κάθε $f \in \Hom_R(M, A \times B)$, αντίστοιχα.
\end{remark}

\begin{exercise}{\rm(\cite[Άσκηση~10.3.13]{DF04})}
  \label{ex:DF_10.3.13}
  Αν $F$ ένα πεπερασμένα παραγόμενο ελεύθερο $R$-πρότυπο, τότε να δείξετε ότι 
  \[
  \Hom_R(F,R) \cong F.
  \] 
\end{exercise}

\begin{proof}[Λύση]
  Έστω $\{e_1, e_2, \dots, e_n\}$ μια βάση του $F$. Για κάθε $1 \le i \le n$, θεωρούμε τα στοιχεία
  \begin{align*}
    e_i^\ast : F &\to R \\
    e_j &\mapsto \delta_{ij} = 
  \begin{cases}
    1, &\ \text{αν $i=j$} \\
    0, &\ \text{διαφορετικά}.
  \end{cases}
  \end{align*}
 του $\Hom_R(F,R)$. Από το Πόρισμα 5.4, αρκεί να δείξουμε ότι το σύνολο $\{e_1^\ast, e_2^\ast, \dots, e_n^\ast\}$ αποτελεί βάση του $\Hom_R(R,F)$.
  \begin{itemize}
    \item Έστω $\phi \in \Hom_R(F,R)$. Τότε 
    \[
    \phi = s_1e_1^\ast + s_2e_2^\ast + \cdots + s_ne_n^\ast
    \] 
    $s_i = \phi(e_i)$, για κάθε $1 \le i \le n$. Πράγματι, αν $a \in F$, τότε 
    \[
    a = r_1e_1 + r_2e_2 + \cdots + r_ne_n
    \]
    για κάποια $r_1, r_2, \dots, r_n \in R$. Συνεπώς,  
    \[
    e_i^\ast(a) = r_i
    \]
    για κάθε $1 \le i \le n$ και γι αυτό 
    \begin{align*}
    \phi(a) 
    &= r_1\phi(e_1) + r_2\phi(e_2) + \cdots + r_n\phi(e_n) \\ 
    &= s_1e_1^\ast(a) + s_2e_2^\ast(a) + \cdots + s_ne_n^\ast(a).
    \end{align*}
    Άρα, το σύνολο $\{e_1^\ast, e_2^\ast, \dots, e_n^\ast\}$ παράγει το $\Hom_R(F,R)$.

    \item Υποθέτουμε ότι 
    \[
    s_1e_1^\ast + s_2e_2^\ast + \cdots + s_ne_n^\ast = 0_{\Hom_R(F,R)}
    \]
    για κάποια $s_1, s_2, \dots, s_n \in R$. Για κάθε $1 \le i \le n$ 
    \begin{align*}
      0_R &= \left(
        s_1e_1^\ast + s_2e_2^\ast + \cdots + s_ne_n^\ast
      \right)(e_i) \\
      &= s_1e_1^\ast(e_i) + s_2e_2^\ast(e_i) + \cdots + s_ne_n^\ast(e_i) \\
      &= s_i.
    \end{align*}
    Επομένως,  $s_i = 0_R$, για κάθε $1 \le i \le n$ και γι αυτό το σύνολο $\{e_1^\ast, e_2^\ast, \dots, e_n^\ast\}$ είναι $R$-γραμμικά ανεξάρτητο.
  \end{itemize}
\end{proof}

\begin{exercise}{\rm(Κινέζικο θέωρημα υπολοίπων)}
  \label{ex:chinese_remainder_thm}
  Έστω $I, J$ δύο ιδεώδη του $R$. Να δείξετε τα εξής.
  \begin{itemize}
    \item[(α)] Υπάρχει ακριβής ακολουθία 
    \[
    \begin{tikzcd}
        0 \arrow[r] & I \cap J \arrow[r] & R \arrow[r] & R/I \times R/J \arrow[r] & R/(I+J) \arrow[r] & 0.
    \end{tikzcd}
    \]
    \item[(β)] Αν $I + J = R$, τότε 
    \[
    R/(I\cap J) \cong R/I \times R/J.
    \]
    \item[(γ)] Αν $n$ και $m$ είναι σχετικά πρώτοι θετικοί ακέραιοι, τότε 
    \[
    \ZZ_{nm} \cong \ZZ_n \times \ZZ_m.
    \]
  \end{itemize}
\end{exercise}

\begin{proof}[Λύση]
    Θα δείξουμε μόνο το (α). Το (β) έπεται από το (α) και το (γ) από το (β). Οι λεπτομέρειες αφήνονται ως άσκηση. Έστω $\iota : I \cap J \to R$ ο \textquote{φυσικός} $R$-μονομορφισμός. Θεωρούμε την απεικόνιση 
    \begin{align*}
      \phi : R &\to R/I \times R/J \\
            r &\mapsto (r + I, r + J).
    \end{align*}
    Αυτή είναι $R$-γραμμική (γιατί;) και $\ker(\phi) = \rmIm(\iota)=I\cap{J}$. Πράγματι,
    \begin{align*}
      r \in \ker(\phi) 
      \ \Leftrightarrow \ \phi(r) = (I, J) 
      \ \Leftrightarrow \ (r+I, r+J) = (I, J) 
      \ \Leftrightarrow \ r \in I\cap J. 
    \end{align*}
    Έχουμε λοιπόν, μία ακριβή ακολουθία 
    \[
    \begin{tikzcd}
        0 \arrow[r] & I \cap J \arrow[r, hook, "\iota"] & R \arrow[r, "\phi"] & R/I \times R/J.
    \end{tikzcd}
    \]

    Έπειτα, θεωρούμε την απεικόνιση 
    \begin{align*}
      \psi: R/I \times R/J &\to R/(I+J) \\
      (r+I, s+J) &\mapsto (r-s) + (I+J).
    \end{align*}
    Η $\psi$ είναι 
    \begin{itemize}
      \item καλώς ορισμένη, δηλαδή 
      \[
      r + I = r' + I, \ s + J = s' + J \quad \then \quad 
      \phi(r+I,s+J) = \phi(r'+I,s'+J).
      \]
      Πράγματι, αν $r + I = r' + I$ και $s + J = s' + J$, τότε $r-r' \in I$ και $s-s' \in J$ και γι αυτό 
      \[
      (r-s) - (r'-s') = (r-r') - (s-s') \ \in \ I+J
      \]
      που σημαίνει ότι 
      \[
      (r-s) + (I+J) = (r'-s') + (I+J) 
      \]
      και γι αυτό $\psi(r+I, s+J) = \psi(r'+I, s'+J)$, όπως θέλαμε.
      \item $R$-ομομορφισμός (αφήνεται ως άσκηση) και 
      \item επί, διότι αν $r \in R$, τότε 
      \[
      \psi(r+I, 0_R + J) = r + (I+J).
      \]
    \end{itemize}

    Τέλος, παρατηρούμε ότι $\ker(\psi) = \rmIm(\phi)$. Πράγματι, αν $(r+I, s+J) \in \ker(\psi)$, τότε $r-s \in I+J$ και γι αυτό 
    \[
    r-s = i + j,
    \]
    για κάποια $i \in I$ και $j \in J$. Συνεπώς, για το στοιχείο 
    \[
    a = r-i = s+j
    \]
    έχουμε  
    \begin{align*}
    \phi(a) &= 
    \left(
      a + I, a + J
    \right) \\
    &= \left(
      r-i + I, s-j + J 
    \right) \\
    &= \left(
      r+I, s+J
    \right)
    \end{align*}
    και γι' αυτό $(r+I, s+J) \in \rmIm(\phi)$. Επομένως, 
    $
    \ker(\psi) \subseteq \rmIm(\phi)
    $. 
    Αντιστρόφως, αν $a \in \rmIm(\phi)$, τότε $a = (r+I, r+J)$ για κάποιο $r \in R$ και γι αυτό 
    \begin{align*}
    \psi(a) 
    &= \psi(r+I,r+J) \\
    &= (r-r) + (I+J) \\ 
    &= I+J, 
    \end{align*}
    που σημαίνει ότι $a \in \ker(\psi)$. Συνεπώς, $\rmIm(\phi) \subseteq \ker(\psi)$.

    Άρα, κατασκευάσαμε μία ακριβή ακολουθία 
    \[
    \begin{tikzcd}
        0 \arrow[r] & I \cap J \arrow[r, hook, "\iota"] & R \arrow[r, "\phi"] & R/I \times R/J \arrow[r, "\psi"] & R/(I+J) \arrow[r] & 0
    \end{tikzcd}
    \]
    (γιατί;).
\end{proof}

\begin{remark}
  Ο ισομορφισμός $R/(I\cap{J}) \cong R/I \times R/J$ ισχύει και στο επίπεδο των δακτυλίων, διότι η απεικόνιση $\phi$ είναι ομομορφισμός δακτυλίων. Πράγματι, αν $r, s \in R$, τότε 
  \[
  \phi(rs) = \left(
    rs + I, rs + J
  \right) = \left(
    r + I, r + J
  \right)\left(
    s + I, s + J
  \right) = \phi(r)\phi(s).
  \]
\end{remark}

\begin{thebibliography}{99}
    \bibitem{DF04}
    D. S. Dummit και R. M. Foote, 
    \textit{Abstract algebra},
    third edition, Wiley, 2004. 
    % \bibitem{Mal08}
    % Μ. Μαλιάκας, 
    % \textit{Εισαγωγή στην Μεταθετική Άλγεβρα}, 
    % Εκδόσεις Σοφία, 2008.
\end{thebibliography}
\end{document}